\documentclass[11pt]{article}
\usepackage[a4paper,margin=25mm]{geometry}
\usepackage{amsmath,amssymb,amsthm}
\usepackage[hidelinks]{hyperref}
\emergencystretch=2em
\newtheorem{theorem}{Theorem}
\title{An open queue with inverse-deficit length scale\newline
\large Disproof of Conjecture 00000005501}
\author{gaochengzhecpu}\date{}
\begin{document}\maketitle
\begin{abstract}
A single stationary M/M/1 queue is an open Markovian
network with exponential service. With service rate one
and load $\rho\in(0,1)$, its mean population is
$\rho/(1-\rho)$, with deficit exponent $-1$.
Moreover, along an explicit stable sequence approaching
load one, every fixed multiple of the proposed
$(1-\rho)^{-1/2}$ scale captures probability tending to
zero. The contradiction therefore concerns the actual
probability scale as well as the mean, and holds whether
queue length includes the customer in service or counts
only waiting customers.
\end{abstract}

\section*{An actual open Markovian queue}
Take an infinite-capacity station with external Poisson
arrivals of rate $\rho\in(0,1)$, a single exponential
server of rate one, and departures out of the network.
This is an open M/M/1 network; the source does not require
more than one station. Its load is $\rho$ and its
stability boundary is $\rho=1$. Write $\delta=1-\rho$.
The standard queue model and its birth/death rates are
described in \cite{modiano}.

Here is an explicit equivalent construction used in Lean.
Run a Poisson clock of rate $1+\rho$. At each clock tick,
independently choose arrival with probability
$u=\rho/(1+\rho)$, and potential service with probability
$d=1/(1+\rho)$. An arrival changes $n$ to $n+1$.
Potential service changes $n$ to $(n-1)^+$; at zero it
is an idle event. The actual transition probability is
\[
 T(n,m)=u\boldsymbol1_{m=n+1}
           +d\boldsymbol1_{m=(n-1)^+}.
\]
Its Poisson uniformization generator $Q=(1+\rho)(T-I)$
has the complete list of entries
\begin{gather*}
 Q(n,n+1)=\rho,\quad Q(n+1,n)=1,\quad Q(0,0)=-\rho,\\
 Q(n+1,n+1)=-(1+\rho),
 \qquad Q(n,m)=0\quad\text{otherwise}.
\end{gather*}
These are precisely the M/M/1 rates. Thus the geometric
distribution used below is tied to a genuine queue
transition law, rather than merely named as a queue length.

\newpage
\section*{Stationarity and the true expectation}
Let $\pi_\rho$ be the normalized geometric probability law
\[
                    \pi_\rho(n)=\delta\rho^n,
                      \qquad n=0,1,\ldots.
\]
Normalization follows from the convergent geometric
series. Adjacent entries satisfy detailed balance since
\[
 \pi_\rho(n)u
   =\delta\rho^n\frac\rho{1+\rho}
   =\delta\rho^{n+1}\frac1{1+\rho}
   =\pi_\rho(n+1)d.
\]
For nonadjacent distinct states both transition terms
are zero, and diagonal detailed balance is automatic.
Summing the identity over the first state gives
$\pi_\rho T=\pi_\rho$. It follows by induction that
$\pi_\rho T^k=\pi_\rho$ for every $k$.

At time $t\geq0$ the actual uniformized queue law,
initialized with $\pi_\rho$, is
\[
 \sum_{k=0}^{\infty}
   e^{-(1+\rho)t}\frac{((1+\rho)t)^k}{k!}
             (\pi_\rho T^k)=\pi_\rho.
\]
This proves all-time stationarity for the explicitly
specified queue. In Lean, $T$ is an actual PMF transition,
iterations are actual PMF binds, and the mixture is
Mathlib's actual Poisson PMF bind.

The expectation is the genuine nonnegative Lebesgue
integral of state $n$ against this probability measure.
The weighted geometric series gives the finite value
\[
 \mathbb E_\rho N
   =\sum_{n=0}^{\infty}n\delta\rho^n
   =\delta\frac\rho{(1-\rho)^2}
   =\frac\rho\delta,
       \qquad \delta\mathbb E_\rho N=\rho.
\]
In particular its deficit exponent is $-1$ and
$\delta\mathbb E_\rho N\to1$ as $\rho\uparrow1$.
This disagrees with a finite nonzero prefactor on the
proposed half-power scale. The next argument also rules
out a rare-tail explanation for this mean divergence.

\section*{An actual distribution bound}
For every $K\in\mathbb N$, each atom has mass at most
$\delta$, because $\rho^n\leq1$. Hence
\[
       0\leq\Pr_\rho(N\leq K)
              =\sum_{n=0}^{K}\delta\rho^n
                    \leq(K+1)\delta.
\]
This is a bound for the actual stationary probability
measure, not a moment inequality or an approximation.

\newpage
\section*{The half-power scale fails in probability}
For $k\geq0$, put
\[
 a_k=k+2,\qquad \rho_k=1-a_k^{-2},\qquad
 \delta_k=a_k^{-2}.
\]
For every $k$, $0<\rho_k<1$, and $\rho_k\to1$.
The proposed scale is exactly
$\delta_k^{-1/2}=a_k$. Its normalized mean is
\[
       \sqrt{\delta_k}\,\mathbb E_{\rho_k}N
              =a_k-a_k^{-1}\longrightarrow+\infty.
\]
For every fixed integer $M\geq0$, the actual cutoff
probability satisfies
\[
 0\leq\Pr_{\rho_k}(N\leq M a_k)
       \leq\frac{M a_k+1}{a_k^2}
       \leq\frac{M+1}{a_k}\longrightarrow0.
\]
Arbitrary finite real cutoffs are covered by choosing
a larger integer $M$ and using event inclusion.
Thus $\sqrt{\delta_k}N$ escapes to $+\infty$ in
probability. In particular, these scaled distributions
are not tight and cannot have a finite stationary
half-power scaling law. The failure along one legitimate
sequence suffices to refute the asserted general
near-boundary scale.

\section*{If queue length means waiting customers only}
Let $W=(N-1)^+$. Its actual distribution is the pushforward
of $\pi_\rho$ by this map. For every integer $K\geq0$,
\[
        \{W\leq K\}=\{N\leq K+1\}.
\]
Consequently
\[
 0\leq\Pr_{\rho_k}(W\leq M a_k)
    \leq\frac{M a_k+2}{a_k^2}
    \leq\frac{M+2}{a_k}\longrightarrow0.
\]
The same probabilistic obstruction therefore holds under
both common conventions for queue length.

\begin{theorem}
An open Markovian network with exponential service need
not have queue length on the negative half-power scale
of the load deficit near stability.
\end{theorem}
\begin{proof}
The single stationary M/M/1 queue above satisfies the
stated network definition, while every finite multiple
of that proposed scale captures probability tending to
zero along $\rho_k\uparrow1$, under either convention.
\end{proof}
The false exponent already refutes the conjunction;
the additional redistribution prefactor need not be
determined. The example concerns stationary queue length.
It makes no claim about centered fluctuations or
fixed-time transient heavy-traffic approximations.

\newpage
\section*{Formalization and reproducibility}
The Lean project defines the normalized arrival/service
coin as an actual PMF, maps it to the actual queue
transition, and proves the displayed matrix entries.
Detailed balance is proved for all state pairs and
summed using PMF row normalization. Iteration and actual
Poisson mixing then prove stationarity at every
nonnegative time. All entries of the uniformization
generator are proved, including the idle-clock boundary.

Poisson uniformization specifies the actual queue model
used in the proof. The formal results concern its
stationary law at every time. A separate construction
of complete paths with exponential holding times, and
a semigroup derivative at zero, are outside the Lean
claims and are unnecessary for this counterexample.

The mean is a nonnegative measure integral, evaluated
using a proved convergent weighted geometric series.
Its conversion to a real value occurs only after proving
the integral equals the finite positive formula; no
infinite-to-real convention enters the calculation.
The cutoff probabilities are actual measure values
on finite sets, with exact finite-sum identities.
The stable approach sequence, its limit and square-root
deficit, the diverging normalized mean, and every
integer-cutoff probability limit are formal theorems.

The waiting-only distribution is an actual pushforward,
and its event identity and vanishing cutoff probabilities
are also formal theorems. Separate statements rule out
eventual capture of even probability $1/2$ by any fixed
integer multiple of the proposed scale. No numerical
or auxiliary verifier is required.

Lean 4.19.0 and Mathlib revision\\
\texttt{c44e0c8ee63ca166450922a373c7409c5d26b00b}
are pinned. From \texttt{lean/}, run
\begin{verbatim}
lake exe cache get
lake build
lake env lean -DwarningAsError=true Main.lean
\end{verbatim}
Fresh-build logs, source and hashes are included.
Printed dependencies contain only standard foundational
axioms. Compile with \texttt{tectonic main.tex}.
The author agent and parent reviewed the work internally;
no external independent review is claimed.

\begin{thebibliography}{9}
\bibitem{source} The Last Math Competition,
Conjecture 00000005501. Exact bilingual source included.
\bibitem{modiano} E. Modiano, MIT 6.263,
\emph{Lectures 5 and 6: Queueing Systems}, M/M/1 queue.
\href{https://web.mit.edu/modiano/www/6.263/lec5-6.pdf}{Author-hosted lecture slides}.
The reference identifies the standard queue model;
the stationarity and probability-scale proof is given above.
\end{thebibliography}
\end{document}
